\documentclass[12pt]{article}

%%% uniform margins
\usepackage[margin=25mm]{geometry}
%%% hyperlinks
\usepackage{hyperref,url}
\usepackage[svgnames]{xcolor}
\hypersetup{
    colorlinks=true,
    linkcolor=DarkBlue,    
    urlcolor=DarkRed,
}

%%% trees
\usepackage{forest}
\useforestlibrary{linguistics}
\forestapplylibrarydefaults{linguistics}

%%% AVMs (Language Science Press)
\usepackage{langsci-avm}
% \avmsetup{
%   font=\scshape,          % attributes in small caps
%   valfont=\itshape,       % values in italics
%   sortfont=\scriptsize\itshape % sorts/types smaller italic
%}

%%% examples
\usepackage{langsci-gb4e}
\let\eachwordone=\itshape

\title{A very short guide to formatting HPSG with LaTeX}
\author{Francis Bond
\\ \texttt{bond@ieee.org}}
\date{February 2026}

\begin{document}
\maketitle

Some recommended packages, and examples of trees and feature structures.  For some advice on style take a look at my \href{https://fcbond.github.io/static/img/ling-style.pdf}{\textit{(Computational) Linguistic Style Guidelines: a guide for the flummoxed}}, or the very comprehensive \href{https://langsci.github.io/gsr/GenericStyleRulesLangsci.pdf}{\textit{Generic Style Rules for Linguistics:
Abridged LangSci version}}.
%\section{Prologue}

\section{Syntactic Tree}
For trees, we use \href{https://ctan.org/pkg/forest}{forest}:

\begin{center}
  
\begin{forest}
for tree={
    if n children=0{font=\itshape, tier=terminal}{},
}
  [S [NP [N [Kim]]]
     [VP [V [seemed]]
         [AP [A [close]] [PP [P [to]] [NP [N [me]]]]] 
         [PP [P [to]] [NP [N [Sandy]]]]]]
\end{forest}
\end{center}

You can draw the trees using brackets.

\section{AVM}
We use \href{https://ctan.org/pkg/langsci-avm}{langsci-avm} for feature structures (attribute value matrices).

\subsection*{Lexical Entry for \textit{eat}}
\begin{center}
\avm{
< \textnormal{devour},  
[ \type*{word} \\ 
  head & \type{verb} \\
  val & [ comps & < NP > ] ] >
}
\end{center}

\subsection*{Binary Head-Complement Rule}

\begin{center}
\begin{tabular}{ccc}
\avm{[ \type*{phrase} \\
   val & [ comps & < > ] ]
   }     
   &  $\to$ & 
  \textbf{H} \avm{[ \type*{word} \\
     val & [ comps & < \1 > ] ]\ \ \1}
\end{tabular}
\end{center}



\section{Tree and AVM}

AVMs can be embedded in tree nodes:

\begin{forest}
  [ S \\ \avm{ [val & [ spr & <> \\ comps & <>]]}
   [ \avm{\1} NP \\ \avm{[val & [ spr & <> \\ comps & <>]]} 
   [ N [ Kim] ] ]
    [ {VP \\ \avm{[val & [ spr & < \1 > \\ comps & <>]]}}
    [ V \\ \avm{[  spr & < \1 > \\comps & < \2 > ]} [ loves] ]
    [\avm{\2} NP  [{every clever linguist}, roof] ]]]
  \end{forest}   

  It's easiest to build the tree first, then add the avms bit by bit, \ldots
  
\section{Type Hierarchy}

You can also add extra lines to trees, using the underlying package \textbf{tikz}.  See the \textbf{forest} documentation for more details.

\begin{forest}
 /tikz/every node/.append style={font=\it},
 [agr-cat [sg,name=sg [3sg]] [,phantom] [,phantom] [,phantom]
          [non-3sg [1sg,name=1sg] [non-1sg [2sg,name=2sg] [pl]]]]
 \draw (sg.south) -- (1sg.north);
 \draw (sg.south) -- (2sg.north);
\end{forest}

\section{Numbered Examples and Glosses}

I recommend using  \href{https://ctan.math.utah.edu/ctan/tex-archive/macros/xetex/latex/langsci/documentation/langsci-gb4e.pdf}{langsci-gb4e} for numbered examples.

\begin{exe}
\ex An example
\ex Dalši příklad “Another example” (Czech) \label{s:dalsi}
\ex Japanese \label{s:mammoth}   \\
    \gll manmosu-wa zetumetu-shita \\
         mammoth-{\sc top} die.out-did \\
    \trans “Mammoths died out.” 
\ex Another example
\end{exe}

You can refer to an example: (\ref{s:dalsi}) is an informal way, only ok if you know your readers will know the languages;  (\ref{s:mammoth}) follows the Leipzig Glossing Rules.

In your assignment, you should gloss them with the \href{https://www.eva.mpg.de/lingua/pdf/Glossing-Rules.pdf}{Leipzig Glossing Rules}.

\end{document}
